\section{引言}

\subsection{实验背景}

量子计算的线路模型给出了一个清晰的抽象：量子寄存器保存量子态，量子门按时间顺序
作用在寄存器上，整个计算过程对应复向量空间中的酉演化。课程中讨论的叠加、纠缠、
测量、CNOT、QFT 和 QAOA 等概念，都可以在线路模型中得到统一描述。这个模型适合
讲清楚量子信息的数学结构，也适合作为算法设计的起点。

真实芯片给这个模型增加了硬件约束。以超导量子芯片为例，物理量子比特按固定拓扑
连接，双比特门只能在相邻物理量子比特之间直接执行。理想线路中的两个逻辑量子比特
可能在算法层面相互作用频繁，但映射到芯片后，它们未必位于相邻物理位置。编译器
需要在逻辑线路和物理拓扑之间建立对应关系，并在不相邻时插入 SWAP 重新安排量子态。
这一步不会改变理想语义，却会增加额外门操作。在线路深度增大后，退相干、门误差和
串扰带来的影响更明显，路由质量会直接影响量子程序在真实设备上的可执行性和可靠性。

量子线路路由因此成为一个典型的交叉问题。一方面，它建立在量子信息基础之上：每个
门都是对量子态的酉变换，SWAP 对应两个量子比特状态的交换，硬件拓扑限制了双比特
相互作用的物理位置。另一方面，它又是一个组合优化问题：编译器必须在大量可能映射
和候选 SWAP 中选择一条路径，使输出线路满足拓扑约束，同时尽量减少 SWAP 数、深度
和编译耗时。对于几十个量子比特的线路，简单枚举不可行，启发式搜索和学习辅助方法
具有实际意义。

\subsection{实验目标}

本项目从复现 SABRE 基线开始。SABRE 是量子路由中常用的启发式方法，它利用前沿门、
物理距离和后续门信息选择候选 SWAP，能够在较短时间内给出稳定结果。先实现并复现
这一基线，有两个作用：第一，它给出了传统方法的可比较参照；第二，它让后续自研
方法始终在同一线路、同一拓扑和同一指标下讨论，而不是只展示单独算法的绝对数值。

在此基础上，项目设计了 NPQR 学习评分路由方法。NPQR 并不把路由完全交给学习模型
决策，而是把模型放入显式搜索流程中：初始映射根据逻辑交互结构选择起点，动作掩码
缩小候选 SWAP 集合，模型评分给候选动作排序，束搜索保留多条可能路线，局部修复
处理尾部卡顿状态，轨迹复放验证每一步门是否满足硬件边约束。这样的设计保留了搜索
过程的可解释性，也使结果能够通过可视化回放和指标表格相互印证。

项目最终不只停留在算法脚本层面，而是实现为一个可演示、可运行、可复现的系统成果。
系统支持网页演示、服务器接口、工具接口和本地命令行；前端展示线路输入、拓扑映射、
SWAP 插入和指标对比；后端负责执行编译、返回结果并支持扩展实验。围绕这些交付形式，
项目补齐 README、测试、部署说明和报告材料，使代码实现、在线页面、服务接口和
本地示例形成统一交付。

\subsection{实验设计思路}

本文按“基础问题、基线复现、自研方法、系统交付、实验验证”的顺序展开。后文先介绍
量子比特、量子门、SWAP、硬件拓扑和路由目标，再说明 SABRE 基线的复现方法与实验
口径；随后给出 NPQR 的自然演进过程，并介绍系统界面、轨迹回放、部署方式和测试
验证；在过程说明部分，报告单独讨论 AI 协作在资料整理、开发、测试和写作中的使用
方法；实验部分集中呈现主实验、扩展实验、消融实验和耗时分析；最后总结课程知识、
算法实现、工程交付和后续优化方向。
